%% =====================================================================
%%  B4-T-24-03 -- what B4 contributes to "Refusing a Deadline"
%%  -------------------------------------------------------------------
%%  LAYER A: APPEND-ONLY. Written once, then left alone. Material found
%%  only here sits in the `unique` slot and is protected by rule R10.
%% =====================================================================
%% @kind: source
%% @id: B4-T-24-03
%% @book: B4
%% @topic: T-24-03
%% @locator: ch. 7, pp. 205--210
%% @status: distilled
%% @unique: yes
%% @integrated: yes
%% @updated: 2026-09-19

\dossier{B4}{T-24-03}{\bookshort{B4}}{ch. 7, pp. 205--210}

\begin{distilled}
  \item Time pressure is the condition under which the other principles in the book operate most effectively, since it removes the opportunity to check any of them.
  \item The recommended response is to address the timing rather than the offer, which keeps the proposal available while removing the condition that makes a poor decision likely.
  \item A genuine constraint produces a consequence when ignored; a manufactured one produces a replacement deadline.
\end{distilled}

\begin{unique}
  \item The clock-versus-offer distinction as the structure of the defence, with the practical advantage that arguing about timing is ground the claimant is weaker on.
  \item The ignore test as the diagnostic that separates the two cases.
\end{unique}
